\documentclass[twocolumn]{article}

\usepackage{amsmath, amssymb}
\usepackage{booktabs}
\usepackage{array}
\usepackage{natbib}
\usepackage{hyperref}
\usepackage{xcolor}
\usepackage{microtype}
\usepackage[margin=0.85in]{geometry}
\usepackage{graphicx}
\usepackage{caption}
\usepackage{enumitem}

\hypersetup{colorlinks=true,linkcolor=blue,citecolor=blue,urlcolor=blue}

\newcolumntype{L}[1]{>{\raggedright\arraybackslash}p{#1}}
\newcolumntype{C}[1]{>{\centering\arraybackslash}p{#1}}

\begin{document}

\title{Cluster Morphology as a Systematic in ISW Measurements:\\
Evidence from Planck SZ Clusters and SZ-Free CMB Maps}

\author{Alika M.\ Parks}

\date{}
\twocolumn[
  \begin{@twocolumnfalse}
  \maketitle

  \vspace{-10pt}
  \noindent\textbf{Alika M.\ Parks}\\
  Independent Researcher, Kalaheo, HI, USA\\
  \href{mailto:alikamp@gmail.com}{alikamp@gmail.com}

  \vspace{12pt}
  \noindent\rule{\textwidth}{0.4pt}

  \vspace{6pt}
  \noindent\textbf{Abstract}\quad
  We test whether the dynamical state of galaxy clusters correlates
  with the Integrated Sachs-Wolfe (ISW) temperature signal measured
  at their locations.
  Using 942 confirmed clusters from the Planck PSZ2 catalogue with
  the Planck SMICA SZ-free CMB map, we classify clusters as relaxed
  or disturbed based on their residual from the $Y_{\mathrm{SZ}}$--mass
  scaling relation.
  Disturbed clusters show a mean CMB temperature decrement of
  $\langle\Delta T\rangle = -5.1 \pm 2.2\;\mu\mathrm{K}$, while
  relaxed clusters show $\langle\Delta T\rangle = +2.0 \pm 2.6\;\mu\mathrm{K}$,
  yielding a morphology-dependent difference of
  $-7.1 \pm 3.4\;\mu\mathrm{K}$ ($2.1\sigma$, $p = 0.037$).
  After mass-matching the two subsamples (364 pairs,
  KS test $p = 0.41$), the difference persists at
  $-6.5 \pm 3.8\;\mu\mathrm{K}$ ($1.7\sigma$, $p = 0.087$),
  indicating the effect is not driven by cluster mass.
  Multivariate regression controlling for mass and redshift yields a
  morphology coefficient of $-16.6 \pm 9.3\;\mu\mathrm{K}$ ($1.8\sigma$).
  The signal is concentrated at $z < 0.2$, with the strongest
  contribution at $0.10 < z < 0.15$ ($-28.7\;\mu\mathrm{K}$,
  $2.8\sigma$, $p = 0.006$).
  The sign and redshift dependence are consistent with recent reports
  of an anomalous negative ISW effect in the nearby Universe.
  We discuss implications for ISW stacking analyses and the
  interpretation of the late-time ISW anomaly.

  \vspace{8pt}
  \noindent\textbf{Keywords}\quad
  Cosmic microwave background $\cdot$
  Integrated Sachs-Wolfe effect $\cdot$
  Galaxy clusters $\cdot$
  Large-scale structure $\cdot$
  Cluster morphology $\cdot$
  Dark energy

  \vspace{4pt}
  \noindent\rule{\textwidth}{0.4pt}
  \vspace{12pt}
  \end{@twocolumnfalse}
]

% ─────────────────────────────────────────────────────────────────────────────
\section{Introduction}
\label{sec:intro}

The late-time Integrated Sachs-Wolfe (ISW) effect arises when CMB
photons traverse time-varying gravitational potentials in a
dark-energy-dominated Universe \citep{sachs1967,crittenden1996}.
In the standard $\Lambda$CDM cosmology, gravitational potentials
decay at late times as dark energy suppresses structure growth,
producing a net blueshift for photons passing through overdensities
and a net redshift for photons crossing voids.
Detection of the ISW effect through cross-correlation of CMB maps
with large-scale structure tracers has been reported at
$2$--$5\sigma$ significance
\citep{granett2008,planckisw2016}.

However, recent measurements have revealed significant anomalies.
\citet{hansen2024} and \citet{hansen2026} report that CMB photons
passing through nearby ($z < 0.03$) voids are significantly
\emph{hotter} than predicted, while photons passing through nearby
overdensities are \emph{colder} --- the opposite of the standard
ISW prediction, at $> 6.5\sigma$ combined significance.
Independently, DESI DR2 BAO measurements show $3.1$--$4.2\sigma$
preference for dynamical dark energy over a cosmological constant
\citep{desi2025}.
These anomalies motivate a reassessment of the assumptions
underlying ISW measurements.

Standard ISW calculations assume that the gravitational potential
$\Phi$ of a structure evolves according to the linear growth factor,
which depends only on the background cosmology.
All structures of the same mass and redshift are treated identically
regardless of their internal dynamical state.
However, galaxy clusters span a wide range of morphologies, from
relaxed, centrally concentrated systems to highly disturbed,
actively merging configurations.
Disturbed clusters have flatter pressure profiles, significant
substructure, and are undergoing rapid structural reorganisation
\citep{lovisari2017}.
Whether this internal dynamical state affects the rate of potential
evolution --- and thus the ISW signal --- has not been tested.

In this paper, we perform the first explicit test of whether cluster
morphology correlates with the ISW temperature signal, using 942
confirmed Planck SZ clusters and the Planck SMICA SZ-free CMB map.
We classify clusters as relaxed or disturbed based on their residual
from the $Y_{\mathrm{SZ}}$--mass scaling relation, control for mass
via matched subsamples and multivariate regression, and examine the
redshift dependence of the signal.

% ─────────────────────────────────────────────────────────────────────────────
\section{Data}
\label{sec:data}

\subsection{Cluster Catalogue}

We use the Planck PSZ2 union catalogue
\citep{planckpsz2}, which contains 1653 Sunyaev-Zeldovich-selected
cluster candidates.
We select confirmed clusters (VALIDATION $\geq 20$) with valid
redshift, mass ($M_{500}$), and integrated Compton parameter
($Y_{5R500}$), and apply a galactic latitude cut $|b| > 15^\circ$
to reduce foreground contamination.
This yields $N = 942$ clusters spanning $0.01 < z < 0.97$ and
$M_{500} = 0.79$--$16.1 \times 10^{14}\,M_\odot$.

\subsection{CMB Map}

We use the Planck 2018 SMICA component-separated CMB temperature
map with the SZ signal explicitly deprojected
(\texttt{COM\_CMB\_IQU-smica-nosz\_2048\_R3.00}), at
$N_{\mathrm{side}} = 2048$.
This map is constructed to minimise thermal SZ contamination at
cluster locations, which is critical for this analysis since the
SZ decrement could mimic a negative ISW signal.

% ─────────────────────────────────────────────────────────────────────────────
\section{Methods}
\label{sec:methods}

\subsection{Morphology Classification}

In the absence of X-ray morphology data for the full PSZ2 sample,
we use the residual from the $Y_{\mathrm{SZ}}$--$M_{500}$ scaling
relation as a morphology proxy.
Disturbed clusters undergoing mergers have flatter, more extended
pressure profiles, producing lower integrated SZ signal for their
mass than relaxed, centrally concentrated clusters
\citep{lovisari2017,planckesz2011}.
We fit $\log Y_{\mathrm{SZ}} = a \log M_{500} + b$ to the full
sample and classify clusters with negative residuals
($Y_{\mathrm{SZ}}$ below the fit) as disturbed.
This yields 522 disturbed and 420 relaxed clusters.

\subsection{Aperture Photometry}

At each cluster location, we measure the CMB temperature using
compensated aperture photometry: the mean temperature within a
$15'$ disk minus the mean temperature in a $15'$--$45'$ annulus.
This filter isolates the local temperature excess or decrement
associated with the cluster potential while subtracting the local
CMB background.

\subsection{Mass Matching}

To control for mass dependence, we construct matched subsamples by
pairing each disturbed cluster with the nearest relaxed cluster in
$M_{500}$ within a 20\% tolerance.
This yields 364 matched pairs with statistically indistinguishable
mass distributions (KS test $p = 0.41$).

\subsection{Statistical Tests}

We apply three complementary tests:
(1) Welch's $t$-test comparing the mean $\Delta T$ of relaxed and
disturbed subsamples;
(2) paired $t$-test on the 364 mass-matched pairs;
(3) multivariate OLS regression of $\Delta T$ on $M_{500}$,
redshift $z$, a binary morphology indicator, and a mass--morphology
interaction term.
All significance estimates are cross-checked with 10,000-resample
bootstrap confidence intervals.

\subsection{LLM Disclosure}

Large language model assistance (Claude, Anthropic) was used for
pipeline development and manuscript preparation.
All scientific decisions, analysis methodology, and interpretation
are solely the work of the authors.

% ─────────────────────────────────────────────────────────────────────────────
\section{Results}
\label{sec:results}

\subsection{Full Sample}

The mean CMB temperature at relaxed cluster locations is
$\langle\Delta T\rangle_{\mathrm{rel}} = +2.0 \pm 2.6\;\mu$K
($n = 420$), while at disturbed cluster locations it is
$\langle\Delta T\rangle_{\mathrm{dis}} = -5.1 \pm 2.2\;\mu$K
($n = 522$).
The difference of $-7.1 \pm 3.4\;\mu$K is significant at $2.1\sigma$
($p = 0.037$).
Bootstrap resampling yields a 95\% confidence interval of
$[-13.9, -0.6]\;\mu$K, excluding zero.

\subsection{Mass-Matched Sample}

After mass matching (364 pairs, KS $p = 0.41$), the difference
persists: $-6.5 \pm 3.8\;\mu$K ($1.7\sigma$, $p = 0.087$;
paired $t$-test: $1.7\sigma$, $p = 0.088$).
The reduction in significance relative to the full sample is
consistent with the loss of statistical power from reduced sample
size rather than removal of a mass-driven effect.
The bootstrap 95\% CI for the mass-matched sample is
$[-14.0, +1.0]\;\mu$K.

\subsection{Regression}

Multivariate regression controlling for $M_{500}$ and $z$
simultaneously yields a morphology coefficient of
$-16.6 \pm 9.3\;\mu$K ($1.8\sigma$, $p = 0.073$).
Neither $M_{500}$ ($0.3\sigma$) nor $z$ ($0.9\sigma$) is
individually significant, indicating that the morphology indicator
captures information not explained by mass or redshift alone.

\begin{table}[h]
\centering
\caption{Regression results: $\Delta T = \beta_0 + \beta_1 M_{500}
+ \beta_2 z + \beta_3\,\mathrm{disturbed} +
\beta_4 (M_{500} \times \mathrm{disturbed})$.}
\label{tab:regression}
\begin{tabular}{lrrrr}
\toprule
Parameter & Est. & SE & $t$ & $p$ \\
\midrule
Intercept & 3.06 & 5.64 & 0.54 & 0.587 \\
$M_{500}$ & 0.41 & 1.22 & 0.34 & 0.736 \\
$z$ & $-16.71$ & 18.40 & $-0.91$ & 0.364 \\
Disturbed & $-16.63$ & 9.28 & $-1.79$ & 0.073 \\
$M_{500} \times$ Dist. & 2.46 & 1.97 & 1.25 & 0.212 \\
\bottomrule
\end{tabular}
\end{table}

\subsection{Redshift Dependence}

Table~\ref{tab:zbins} presents the morphology--ISW difference in
redshift bins.
The signal is concentrated at $z < 0.2$, with the strongest
individual bin at $0.10 < z < 0.15$: $\Delta T_{\mathrm{dis}} -
\Delta T_{\mathrm{rel}} = -28.7\;\mu$K at $2.8\sigma$ ($p = 0.006$).

\begin{table}[h]
\centering
\caption{Morphology--ISW difference by redshift.}
\label{tab:zbins}
\begin{tabular}{lrrrrl}
\toprule
$z$ range & $n_r$ & $n_d$ & $\Delta$ ($\mu$K) & $\sigma$ & \\
\midrule
$[0.05, 0.10)$ & 107 & 44 & $-10.0$ & 1.2 & \\
$[0.10, 0.15)$ & 51 & 45 & $-28.7$ & 2.8 & $**$ \\
$[0.15, 0.20)$ & 57 & 65 & $-8.8$ & 0.9 & \\
$[0.20, 0.30)$ & 82 & 141 & $+5.3$ & 0.7 & \\
$[0.30, 0.50)$ & 73 & 157 & $-4.4$ & 0.6 & \\
\bottomrule
\end{tabular}
\end{table}

The redshift concentration of the signal at $z < 0.2$ is notable
because it coincides with the redshift range where
\citet{hansen2026} report the anomalous negative ISW effect.
At $z > 0.2$, the morphology--ISW difference is consistent with
zero.

% ─────────────────────────────────────────────────────────────────────────────
\section{Discussion}
\label{sec:discussion}

\subsection{Interpretation}

The principal finding is that cluster dynamical state correlates
with the CMB temperature measured at cluster locations, even after
controlling for cluster mass and using an SZ-deprojected CMB map.
Disturbed clusters --- those with SZ signal below the $Y$--$M$
scaling relation, indicative of merging or extended pressure
profiles --- show systematically colder CMB temperatures than
relaxed clusters.

The sign of the effect is striking: standard ISW predicts
\emph{positive} $\Delta T$ at overdensities (potentials decaying
$\to$ net blueshift), yet disturbed clusters show \emph{negative}
$\Delta T$.
This sign reversal is consistent with the hypothesis that
the gravitational potentials of dynamically disturbed structures
evolve differently from relaxed systems --- either decaying more
rapidly, or in extreme cases, temporarily \emph{deepening} as
substructure merges and the potential well reorganises.

This is consistent with the results of \citet{hansen2026}, who
find cold overdensities at $z < 0.03$ and interpret their results
as evidence for a negative ISW effect in the nearby Universe.
Our analysis extends the redshift range to $z \sim 0.15$ and
provides a physical mechanism: the morphological state of the
overdensity modulates the ISW signal.

\subsection{Connection to Dynamical Dark Energy}

DESI DR2 BAO measurements prefer dynamical dark energy at
$3.1$--$4.2\sigma$ \citep{desi2025}.
If the dark energy equation of state evolves with time,
the growth factor and hence the rate of potential decay also
evolve, producing a modified ISW signal.
Our results suggest an additional, complementary effect: even
within a fixed cosmological model, the \emph{internal} dynamical
state of structures modulates the ISW signal.
A complete treatment of the late-time ISW may need to account
for both cosmological evolution and structural morphology.

\subsection{Relation to the Structural-Geometric Framework}

This analysis was motivated by the Parks Node Ejection Protocol
(PNEP), a geometric diagnostic for gravitational hierarchy decay
in few-body stellar systems \citep{parks2026}.
PNEP demonstrates that the internal geometry of a gravitational
system --- specifically, the distance hierarchy among bodies
sampled at dynamically privileged moments --- encodes information
about the system's stability that scalar measures (total mass,
energy) miss.
The present work tests the cosmological analog: whether the
internal geometry (morphology) of a galaxy cluster encodes
information about its potential evolution that mass alone misses.
The affirmative result, while at modest significance, supports the
broader principle that structural geometry matters for
gravitational potential evolution across scales.

\subsection{Limitations}

The $Y$--$M$ residual morphology proxy, while physically
motivated, is an imperfect indicator of dynamical state.
X-ray morphology parameters (concentration, centroid shift)
from dedicated observations \citep{lovisari2017} would provide
a more direct classification.
Our analysis uses such parameters for only ${\sim}120$ clusters;
extending this to the full PSZ2 sample would substantially
increase statistical power.

The mass-matched result ($1.7\sigma$) does not reach conventional
significance.
With 364 pairs and ${\sim}50\;\mu$K noise per cluster, we are at
the statistical limit of the current sample.
Forthcoming cluster catalogues from SPT-3G, AdvACT, and the
Simons Observatory will increase sample sizes by an order of
magnitude.

The aperture photometry approach is standard but coarse.
Matched-filter or optimal-filter techniques adapted to the cluster
angular profile would improve signal extraction.

We have not modelled the theoretical ISW signal from $N$-body
simulations for comparison.
A full treatment would predict $\Delta T_{\mathrm{ISW}}$ for
relaxed and disturbed clusters from hydrodynamical simulations
(e.g., cosmoOWLS, BAHAMAS) and compare with observations.

% ─────────────────────────────────────────────────────────────────────────────
\section{Conclusions}
\label{sec:conclusions}

We report the first test of morphology-dependent ISW signal from
galaxy clusters using Planck data.
Our principal findings are:

\begin{enumerate}[nosep]
  \item Disturbed clusters show a mean CMB temperature
        $7.1 \pm 3.4\;\mu$K colder than relaxed clusters
        ($2.1\sigma$, $p = 0.037$) in the SZ-free Planck CMB map.

  \item The signal survives mass matching ($-6.5\;\mu$K at
        $1.7\sigma$, 364 pairs) and multivariate regression
        controlling for mass and redshift ($-16.6\;\mu$K at
        $1.8\sigma$), indicating it is not driven by cluster mass.

  \item The signal is concentrated at $z < 0.2$, with the
        strongest contribution at $0.10 < z < 0.15$ ($-28.7\;\mu$K,
        $2.8\sigma$, $p = 0.006$).

  \item The sign (cold disturbed overdensities) is consistent with
        the anomalous negative ISW reported by \citet{hansen2026}
        at $z < 0.03$.
\end{enumerate}

While the mass-matched significance ($1.7\sigma$) does not reach
the threshold for a definitive detection, the consistency of the
signal across multiple analysis methods, its concentration at
low redshift, and its concordance with independent ISW anomaly
measurements motivate further investigation with larger cluster
samples and dedicated X-ray morphology data.

All analysis code is publicly available at
\url{https://github.com/alikamp/Parks-Node-Ejection-Protocol}.

% ─────────────────────────────────────────────────────────────────────────────
\section*{Statements and Declarations}

\paragraph{Competing Interests.}
The authors declare no competing interests.

\paragraph{Data Availability.}
This work uses publicly available data from the Planck Legacy
Archive and the Planck PSZ2 catalogue.
Analysis code is available at
\url{https://github.com/alikamp/Parks-Node-Ejection-Protocol}.

% ─────────────────────────────────────────────────────────────────────────────
\bibliographystyle{aasjournal}
\begin{thebibliography}{99}

\bibitem[Crittenden \& Turok(1996)]{crittenden1996}
Crittenden, R.~G., \& Turok, N.\ 1996, Physical Review Letters, 76, 575

\bibitem[DESI Collaboration(2025)]{desi2025}
DESI Collaboration 2025, arXiv:2503.14738

\bibitem[Granett et al.(2008)]{granett2008}
Granett, B.~R., Neyrinck, M.~C., \& Szapudi, I.\ 2008, ApJ, 683, L99

\bibitem[Hansen et al.(2024)]{hansen2024}
Hansen, F.~K., et al.\ 2024, A\&A, 681, L7

\bibitem[Hansen et al.(2026)]{hansen2026}
Hansen, F.~K., et al.\ 2026, arXiv:2506.08832

\bibitem[Lovisari et al.(2017)]{lovisari2017}
Lovisari, L., et al.\ 2017, ApJ, 846, 51

\bibitem[Naidoo(2023)]{naidoo2023}
Naidoo, K.\ 2023, arXiv:2308.13617

\bibitem[Parks(2026)]{parks2026}
Parks, A.~M.\ 2026, submitted to CMDA

\bibitem[Planck Collaboration(2011)]{planckesz2011}
Planck Collaboration 2011, A\&A, 536, A8

\bibitem[Planck Collaboration(2016a)]{planckpsz2}
Planck Collaboration 2016a, A\&A, 594, A27

\bibitem[Planck Collaboration(2016b)]{planckisw2016}
Planck Collaboration 2016b, A\&A, 594, A21

\bibitem[Sachs \& Wolfe(1967)]{sachs1967}
Sachs, R.~K., \& Wolfe, A.~M.\ 1967, ApJ, 147, 73

\end{thebibliography}

\end{document}
